\section{Number Theory}

\begin{notationtable}
$a\mid b$ & $a$ divides $b$. \\
$a\nmid b$ & $a$ does not divide $b$. \\
$\gcd(a_1, \dots, a_n)$ & Greatest common divisor \\
$\lcmop(a_1, \dots, a_n)$ & Least common multiple \\
$p^k\parallel n$ & $p^k$ exactly divides $n$: $p^k\mid n$ and $p^{k+1}\nmid n$. \\
$v_p(n), \ \nu_p(n)$ & $p$-adic valuation of $n$ \\
$a\equiv b\pmod n$ & $a$ is congruent to $b$ modulo $n$. \\
$[a]_n, \ \overline a$ & Residue class of $a$ modulo $n$; $\overline a$ when the modulus is clear \\
$a\bmod n$ & Least nonnegative remainder of $a$ modulo $n$, for $n>0$ \\
$\Z/n\Z$ & Ring of residue classes modulo $n$ \\
$(\Z/n\Z)^\times$ & Multiplicative group of units modulo $n$ \\
$a^{-1}\pmod n$ & Multiplicative inverse of $a$ modulo $n$, for $\gcd(a, n)=1$ \\
$\ord_n(a)$ & Multiplicative order of $a$ modulo $n$, for $\gcd(a, n)=1$ \\
$\mathbb F_q, \ \operatorname{GF}(q)$ & Finite field with $q$ elements, where $q$ is a prime power \\
$\displaystyle\left(\frac ap\right)$ & Legendre symbol, for an odd prime $p$ \\
$\displaystyle\left(\frac an\right)$ & Jacobi symbol, for an odd positive integer $n$ \\
$\chi\pmod q$ & Dirichlet character modulo $q$ \\
$\varphi(n)$ & Euler totient function \\
$\lambda(n)$ & Carmichael function: exponent of the finite abelian group $(\Z/n\Z)^\times$; the least positive $m$ such that $a^m\equiv1\pmod n$ for all $a$ coprime to $n$ \\
$B_n$ & $n$th Bernoulli number, under the convention $B_1=-\tfrac12$ \\
$\displaystyle\sum_{d\mid n}$ & Sum over the positive divisors $d$ of $n$ \\
$\displaystyle\prod_{p\mid n}$ & Product over the prime divisors $p$ of $n$ \\
$\displaystyle\sigma_k(n)=\sum_{d\mid n}d^k$ & Sum of the $k$th powers of the positive divisors of $n$ \\
$\tau(n)=\sigma_0(n)$ & Number of positive divisors of $n$ \\
$\sigma(n)=\sigma_1(n)$ & Sum of positive divisors of $n$ \\
$\mu(n)$ & M\"obius function \\
$\omega(n)$ & Number of distinct prime divisors of $n$ \\
$\Omega(n)$ & Number of prime factors of $n$, counted with multiplicity \\
$\rad(n)$ & Product of the distinct prime divisors of $n$ \\
$\displaystyle(f\ast g)(n)=\sum_{d\mid n}f(d)g(n/d)$ & Dirichlet convolution of arithmetic functions $f$ and $g$ \\
$\pi(x)$ & Number of primes not exceeding $x$ \\
$\Lambda(n)$ & Von Mangoldt function: $\ln p$ if $n=p^k$ for a prime $p$ and an integer $k\ge1$, and $0$ otherwise \\
$L(s, \chi)$ & Dirichlet $L$-function, $\displaystyle L(s, \chi)=\sum_{n=1}^{\infty}\frac{\chi(n)}{n^s}$ for $\Re(s)>1$ \\
$[a_0;a_1, a_2, \dots]$ & Simple continued fraction \\
$\displaystyle[a_0;a_1, \dots, a_n]=\frac{p_n}{q_n}$ & $n$th convergent of a continued fraction \\
$\Phi_n(x)$ & $n$th cyclotomic polynomial \\
$\Z[i]$ & Ring of Gaussian integers \\
$K=\Q(\sqrt d)$ & Quadratic field, with squarefree $d\in\Z$ and $d\ne1$ \\
$\mathcal O_K$ & Ring of integers of the number field $K$ \\
$\IdealNorm(\mathfrak a)=|\mathcal O_K/\mathfrak a|$ & Absolute norm of a nonzero integral ideal $\mathfrak a\subseteq\mathcal O_K$ \\
$D_K, \ \disc(K)$ & Discriminant of the number field $K$ \\
$\Cl(K)$ & Ideal class group of the number field $K$ \\
$h_K=|\Cl(K)|$ & Class number of the number field $K$ \\
\end{notationtable}
